\documentclass[11pt,a4paper]{article}
\usepackage[a4paper,margin=1in,headsep=30pt,footskip=35pt]{geometry}
\usepackage[T1]{fontenc}
\usepackage[utf8]{inputenc}
\usepackage{tgpagella}
\usepackage[scale=0.92]{tgheros}
\usepackage{microtype}
\usepackage{graphicx}
\usepackage{hyperref}
\usepackage{enumitem}
\usepackage{xcolor}
\usepackage{titlesec}
\usepackage{parskip}
\usepackage{longtable}
\usepackage{booktabs}
\usepackage{array}
\usepackage{tabularx}
\usepackage{fancyhdr}
\usepackage{lastpage}
\usepackage{tikz}
\usepackage[most]{tcolorbox}

% --- CORE COLOR IDENTITY ---
\definecolor{KazePrimary}{HTML}{284F58}
\definecolor{KazeSecondary}{HTML}{8A6538}
\definecolor{KazeInk}{HTML}{1F292D}
\definecolor{KazeMuted}{HTML}{5D676C}
\definecolor{KazeSoft}{HTML}{F8F4ED}
\definecolor{KazeLine}{HTML}{D8CCBC}
\definecolor{KazePaper}{HTML}{FCFBF8}
\definecolor{KazeLegal}{HTML}{7A2E2E}

\hypersetup{colorlinks=true, linkcolor=KazePrimary, urlcolor=KazePrimary}

\setlist[itemize]{leftmargin=1.5em, itemsep=0.4em}
\renewcommand{\arraystretch}{1.25}
\setlength{\headheight}{35pt}

\newcommand{\KazeLogoPath}{../../composeApp/src/commonMain/composeResources/drawable/k_logo_raster.png}
\newcommand{\KazeMarkPath}{../../composeApp/src/commonMain/composeResources/drawable/k_mark_raster.png}
\newcommand{\GaboLogoPath}{../../composeApp/src/commonMain/composeResources/drawable/gabo_mark_raster.png}
\newcommand{\GaboGrayMarkPath}{../../composeApp/src/commonMain/composeResources/drawable/gabo_mark_gray_raster.png}
\newcommand{\KazeWordmarkHeader}{{\fontfamily{qhv}\selectfont\sffamily\bfseries\small\textcolor{KazePrimary}{KAZE}}}
\newcommand{\KazeWordmarkCover}{{\fontfamily{qhv}\selectfont\sffamily\bfseries\fontsize{30}{30}\selectfont\textcolor{KazePrimary}{KAZE}}}
\newcommand{\KazeByGaboHeader}{{\raisebox{0.05cm}{\includegraphics[height=0.18cm]{\GaboLogoPath}}\hspace{0.2em}\sffamily\tiny\textcolor{KazeSecondary}{by GABO}}}
\newcommand{\KazeByGaboCover}{{\includegraphics[height=0.45cm]{\GaboLogoPath}\hspace{0.4em}\sffamily\large\textcolor{KazeSecondary}{by GABO}}}
\newcommand{\KazeCoverLabel}[1]{{\sffamily\normalsize\textcolor{KazeSecondary}{\textbf{#1}}}}
\newcommand{\KazeCoverTitle}[1]{{\fontfamily{qpl}\selectfont\fontsize{38}{43}\selectfont\textcolor{KazeInk}{#1}}}
\newcommand{\KazeCoverSubtitle}[1]{{\sffamily\Large\textcolor{KazePrimary}{#1}}}
\newcommand{\KazeContactEmail}{orestegabo@icloud.com}
\newcommand{\KazeContactWebsiteUrl}{https://kazerwanda.com}
\newcommand{\KazeContactWebsiteLabel}{kazerwanda.com}
\newcommand{\KazeContactBlock}{
    \begin{itemize}
        \item \textbf{Email:} \href{mailto:\KazeContactEmail}{\KazeContactEmail}
        \item \textbf{Website:} \href{\KazeContactWebsiteUrl}{\KazeContactWebsiteLabel}
    \end{itemize}
}

\pagecolor{KazePaper}

\AddToHook{shipout/background}{%
    \begin{tikzpicture}[remember picture,overlay]
        \fill[KazePrimary, opacity=0.018, rounded corners=14pt]
            ([xshift=0.04\paperwidth,yshift=-0.07\paperheight]current page.north west)
            rectangle
            ([xshift=0.96\paperwidth,yshift=-0.31\paperheight]current page.north west);
        \fill[KazeSecondary, opacity=0.014, rounded corners=14pt]
            ([xshift=0.05\paperwidth,yshift=0.34\paperheight]current page.north west)
            rectangle
            ([xshift=0.95\paperwidth,yshift=0.14\paperheight]current page.south west);

        \draw[KazePrimary, line width=1.7pt, opacity=0.09, line cap=round]
            ([xshift=0.08\paperwidth,yshift=-0.10\paperheight]current page.north west) --
            ([xshift=0.72\paperwidth,yshift=-0.10\paperheight]current page.north west);
        \draw[KazePrimary, line width=0.95pt, opacity=0.085, line cap=round]
            ([xshift=0.12\paperwidth,yshift=-0.135\paperheight]current page.north west) --
            ([xshift=0.88\paperwidth,yshift=-0.135\paperheight]current page.north west);
        \draw[KazeSecondary, line width=1.2pt, opacity=0.08, line cap=round]
            ([xshift=0.18\paperwidth,yshift=0.26\paperheight]current page.south west) --
            ([xshift=0.84\paperwidth,yshift=0.26\paperheight]current page.south west);
        \draw[KazePrimary, line width=1.7pt, opacity=0.09, line cap=round]
            ([xshift=0.14\paperwidth,yshift=0.21\paperheight]current page.south west) --
            ([xshift=0.62\paperwidth,yshift=0.21\paperheight]current page.south west);

        \draw[KazePrimary, line width=1.15pt, opacity=0.09]
            ([xshift=0.62\paperwidth,yshift=-0.06\paperheight]current page.north west) --
            ([xshift=0.76\paperwidth,yshift=-0.06\paperheight]current page.north west) --
            ([xshift=0.82\paperwidth,yshift=-0.11\paperheight]current page.north west) --
            ([xshift=0.93\paperwidth,yshift=-0.11\paperheight]current page.north west) --
            ([xshift=0.93\paperwidth,yshift=-0.18\paperheight]current page.north west) --
            ([xshift=0.82\paperwidth,yshift=-0.18\paperheight]current page.north west) --
            ([xshift=0.75\paperwidth,yshift=-0.24\paperheight]current page.north west) --
            ([xshift=0.58\paperwidth,yshift=-0.24\paperheight]current page.north west);

        \draw[KazeSecondary, line width=1.15pt, opacity=0.08]
            ([xshift=0.10\paperwidth,yshift=0.12\paperheight]current page.south west) --
            ([xshift=0.26\paperwidth,yshift=0.12\paperheight]current page.south west) --
            ([xshift=0.33\paperwidth,yshift=0.17\paperheight]current page.south west) --
            ([xshift=0.48\paperwidth,yshift=0.17\paperheight]current page.south west) --
            ([xshift=0.48\paperwidth,yshift=0.10\paperheight]current page.south west) --
            ([xshift=0.34\paperwidth,yshift=0.10\paperheight]current page.south west) --
            ([xshift=0.27\paperwidth,yshift=0.05\paperheight]current page.south west) --
            ([xshift=0.12\paperwidth,yshift=0.05\paperheight]current page.south west);

        \draw[KazePrimary, line width=1.8pt, opacity=0.05]
            ([xshift=0.70\paperwidth,yshift=-0.22\paperheight]current page.north west) circle [radius=2.0cm];
        \draw[KazeSecondary, line width=1.8pt, opacity=0.048]
            ([xshift=0.14\paperwidth,yshift=0.18\paperheight]current page.south west) circle [radius=2.35cm];

        \draw[KazePrimary, line cap=round, opacity=0.065, line width=1.2pt]
            ([xshift=0.11\paperwidth,yshift=-0.185\paperheight]current page.north west) --
            ([xshift=0.36\paperwidth,yshift=-0.185\paperheight]current page.north west);
        \draw[KazeSecondary, line cap=round, opacity=0.06, line width=0.8pt]
            ([xshift=0.11\paperwidth,yshift=-0.203\paperheight]current page.north west) --
            ([xshift=0.42\paperwidth,yshift=-0.203\paperheight]current page.north west);
        \draw[KazePrimary, line cap=round, opacity=0.065, line width=1.2pt]
            ([xshift=0.11\paperwidth,yshift=-0.221\paperheight]current page.north west) --
            ([xshift=0.48\paperwidth,yshift=-0.221\paperheight]current page.north west);
        \draw[KazeSecondary, line cap=round, opacity=0.06, line width=0.8pt]
            ([xshift=0.11\paperwidth,yshift=-0.239\paperheight]current page.north west) --
            ([xshift=0.54\paperwidth,yshift=-0.239\paperheight]current page.north west);
        \draw[KazePrimary, line cap=round, opacity=0.065, line width=1.2pt]
            ([xshift=0.11\paperwidth,yshift=-0.257\paperheight]current page.north west) --
            ([xshift=0.60\paperwidth,yshift=-0.257\paperheight]current page.north west);
        \draw[KazeSecondary, line cap=round, opacity=0.06, line width=0.8pt]
            ([xshift=0.11\paperwidth,yshift=-0.275\paperheight]current page.north west) --
            ([xshift=0.66\paperwidth,yshift=-0.275\paperheight]current page.north west);
        \draw[KazePrimary, line cap=round, opacity=0.065, line width=1.2pt]
            ([xshift=0.11\paperwidth,yshift=-0.293\paperheight]current page.north west) --
            ([xshift=0.72\paperwidth,yshift=-0.293\paperheight]current page.north west);
        \node[anchor=south east, opacity=0.035] at ([xshift=-0.45cm,yshift=0.65cm]current page.south east) {%
            \includegraphics[width=4.8cm]{\KazeMarkPath}%
        };
        \node[anchor=north east, opacity=0.028] at ([xshift=-1.1cm,yshift=-1.2cm]current page.north east) {%
            \includegraphics[width=2.2cm]{\GaboGrayMarkPath}%
        };
    \end{tikzpicture}%
}

% --- SECTION DESIGN: TECHNICAL PRECISION ---
\titleformat{\section}
{\sffamily\Large\bfseries\color{KazeInk}}
{\thesection}
{1em}
{\MakeUppercase}
[\vspace{0.2em}{\color{KazeSecondary}\hrule height 1.5pt width 2cm}\vspace{0.5em}]

\titleformat{\subsection}
{\sffamily\large\bfseries\color{KazePrimary}}
{\thesubsection}
{1em}
{}

% --- DUAL-BRANDED HEADER (HIGH VISIBILITY) ---
\pagestyle{fancy}
\fancyhf{}
\renewcommand{\headrulewidth}{0pt}
\renewcommand{\footrulewidth}{0pt}

\fancyhead[L]{%
    \begin{tabular}{@{}l l@{}}
        \raisebox{-0.3\height}{\includegraphics[height=0.6cm]{\KazeLogoPath}} &
        \begin{minipage}[b]{4cm}
            \KazeWordmarkHeader\\
            \vspace{-1mm}
            \KazeByGaboHeader
        \end{minipage}
    \end{tabular}
}
\fancyhead[R]{\sffamily\tiny\textcolor{KazeMuted}{KAZE PROPRIETARY // \today}}

\fancyfoot[L]{%
    \begin{tabular}{@{}l@{}}
        {\color{KazeLine}\rule{3.1cm}{0.6pt}}\\[-0.2em]
        {\sffamily\tiny\textcolor{KazeLegal}{CONFIDENTIAL PRIVATE DOCUMENT}}
    \end{tabular}
}
\fancyfoot[C]{%
    \sffamily\scriptsize
    \textcolor{KazeMuted}{\textsc{Page} \thepage\ of \pageref*{LastPage}}
}
\fancyfoot[R]{%
    \begin{tabular}{@{}r@{}}
        {\color{KazeLine}\rule{2.1cm}{0.6pt}}\\[-0.2em]
        {\sffamily\tiny\textcolor{KazePrimary}{Kaze by GABO}}
    \end{tabular}
}

% --- CLEAN REGISTRY BOXES ---
\newtcolorbox{kazehighlightbox}{
    enhanced,
    colback=KazeSoft,      % Only a very faint paper-like tint
    frame hidden,          % No borders at all
    sharp corners,
    left=5mm, right=5mm, top=4mm, bottom=4mm,
% Use a small square bullet as an anchor instead of a line
    overlay={
        \fill[KazeSecondary] ([xshift=2mm,yshift=-4mm]frame.north west) rectangle ([xshift=3mm,yshift=-5mm]frame.north west);
    }
}
\newtcolorbox{kazecalloutbox}[1][]{
    enhanced,
    colback=white,
    colframe=KazeInk,
    boxrule=1pt,
    sharp corners,
    fonttitle=\sffamily\bfseries\small,
    coltitle=white,
    attach boxed title to top left={yshift=-2mm, xshift=5mm},
    boxed title style={colback=KazeInk, sharp corners, boxrule=0pt},
    title=#1,
    left=5mm, top=6mm, bottom=4mm
}

\newcommand{\KazePageTitle}[1]{%
    \vspace{1em}
    {\sffamily\small\textcolor{KazeSecondary}{\textbf{\MakeUppercase{Document Scope: #1}}}\par}
    \vspace{0.5em}
}

\newcommand{\KazeDocumentNotice}{%
    \begin{kazecalloutbox}[Security Disclosure]
    \sffamily\small This dossier contains trade secrets and sensitive intellectual property. \textbf{Unauthorized distribution is strictly prohibited.} Access is limited to Kaze personnel and authorized external stakeholders for legal and investor review.
    \end{kazecalloutbox}
}

% --- THE "Blueprints" COVER PAGE ---
\newcommand{\KazeCover}[4]{
    \hypersetup{pageanchor=false}
    \begin{titlepage}
    \thispagestyle{empty}

% Background "Pillar" for institutional weight
    \begin{tikzpicture}[remember picture,overlay]
    \fill[KazePrimary] (current page.north west) rectangle ([xshift=0.5cm]current page.south west);
    \draw[KazePrimary!55, line width=1.8pt, opacity=0.22, line cap=round]
        ([xshift=0.08\paperwidth,yshift=-0.10\paperheight]current page.north west) --
        ([xshift=0.72\paperwidth,yshift=-0.10\paperheight]current page.north west);
    \draw[KazeSecondary!70, line width=1.1pt, opacity=0.19, line cap=round]
        ([xshift=0.12\paperwidth,yshift=-0.135\paperheight]current page.north west) --
        ([xshift=0.88\paperwidth,yshift=-0.135\paperheight]current page.north west);
    \draw[KazePrimary!55, line width=1.2pt, opacity=0.20]
        ([xshift=0.62\paperwidth,yshift=-0.06\paperheight]current page.north west) --
        ([xshift=0.76\paperwidth,yshift=-0.06\paperheight]current page.north west) --
        ([xshift=0.82\paperwidth,yshift=-0.11\paperheight]current page.north west) --
        ([xshift=0.93\paperwidth,yshift=-0.11\paperheight]current page.north west) --
        ([xshift=0.93\paperwidth,yshift=-0.18\paperheight]current page.north west) --
        ([xshift=0.82\paperwidth,yshift=-0.18\paperheight]current page.north west) --
        ([xshift=0.75\paperwidth,yshift=-0.24\paperheight]current page.north west) --
        ([xshift=0.58\paperwidth,yshift=-0.24\paperheight]current page.north west);
    \draw[KazeSecondary!75, line width=1.35pt, opacity=0.18, line cap=round]
        ([xshift=0.18\paperwidth,yshift=0.26\paperheight]current page.south west) --
        ([xshift=0.84\paperwidth,yshift=0.26\paperheight]current page.south west);
    \draw[KazePrimary!55, line width=1.8pt, opacity=0.20, line cap=round]
        ([xshift=0.14\paperwidth,yshift=0.21\paperheight]current page.south west) --
        ([xshift=0.62\paperwidth,yshift=0.21\paperheight]current page.south west);
    \end{tikzpicture}

    \vspace*{0.5cm}
    \begin{flushleft}
    \hspace{1cm}
    \begin{tabular}{@{}l l@{}}
    \raisebox{-0.3\height}{\includegraphics[width=2.2cm]{\KazeLogoPath}} &
    \begin{minipage}[b]{8cm}
    {\KazeWordmarkCover\par}
    \vspace{0.1cm}
    {\KazeByGaboCover\par}
    \end{minipage}
    \end{tabular}
    \end{flushleft}

    \vspace{3.5cm}

    \hspace{1cm}
    \begin{minipage}{0.85\textwidth}
    \raggedright
    {\KazeCoverLabel{OFFICIAL DOSSIER / #4}\par}
    \vspace{0.95cm}
    {\KazeCoverTitle{#1}\par}
    \vspace{0.7cm}
    {\KazeCoverSubtitle{#2}\par}

    \vspace{1.5cm}
    {\color{KazeLine}\hrule height 2pt width 3cm}
    \vspace{1.5cm}

    {\sffamily\large\textcolor{KazeMuted}{#3}\par}
    \end{minipage}

    \vfill

    \hspace{1cm}
    \begin{minipage}{0.85\textwidth}
    \begin{kazehighlightbox}
    \sffamily\small
    \textbf{\textcolor{KazeLegal}{RESTRICTED DOCUMENT:}} This information is intended for the recipient only. It contains proprietary strategic planning belonging to Kaze. By accessing this document, you agree to non-disclosure terms.
    \end{kazehighlightbox}
    \end{minipage}

    \vspace{1.5cm}
    \hspace{1cm}
    \begin{minipage}{0.85\textwidth}
    \sffamily\tiny\textcolor{KazeMuted}{AUTHORIZATION \& CONTACT:} \\
    \sffamily\small\textbf{\href{mailto:\KazeContactEmail}{\KazeContactEmail}} \hfill \textbf{\MakeUppercase{\KazeContactWebsiteLabel}}
    \end{minipage}

    \end{titlepage}
    \hypersetup{pageanchor=true}
}

\usepackage{listings}
\usepackage{upquote}

\definecolor{CodeBg}{HTML}{F4F1EA}
\definecolor{CodeRule}{HTML}{D8CCBC}
\definecolor{CodeComment}{HTML}{6A737D}
\definecolor{CodeKeyword}{HTML}{284F58}
\definecolor{CodeString}{HTML}{8A6538}

\lstdefinestyle{kazekotlin}{
    language={},
    basicstyle=\ttfamily\small,
    keywordstyle=\color{CodeKeyword}\bfseries,
    commentstyle=\color{CodeComment},
    stringstyle=\color{CodeString},
    showstringspaces=false,
    columns=fullflexible,
    keepspaces=true,
    breaklines=true,
    frame=single,
    rulecolor=\color{CodeRule},
    backgroundcolor=\color{CodeBg},
    xleftmargin=0.2cm,
    xrightmargin=0.2cm,
    aboveskip=0.8em,
    belowskip=0.8em
}

\begin{document}

\KazeCover
  {Ktor Course Notes for Kaze}
  {A beginner-first course for learning APIs and Ktor from scratch}
  {April 2026}
  {Developer learning workbook}

\KazePageTitle{Ktor Course Notes}
\KazeDocumentNotice

\section{Before Chapter 1: What An API Is}

Because you said you have coded in C++ and Java but have never built an API, we should begin one step before Ktor.

\subsection{What Is An API?}

An API is a contract that lets one program talk to another program.

When your mobile app asks a server:
\begin{itemize}
    \item ``please sign this user in''
    \item ``please give me the list of hotels''
    \item ``please store this late checkout request''
\end{itemize}

that conversation is happening through an API.

In practice, most modern app backends expose an HTTP API. That means the client sends HTTP requests such as:
\begin{itemize}
    \item \texttt{GET /hotels}
    \item \texttt{POST /auth/signin}
    \item \texttt{GET /auth/me}
\end{itemize}

and the server sends HTTP responses back, usually in JSON.

\subsection{A Good Beginner Analogy}

Think of:
\begin{itemize}
    \item the mobile app as a customer
    \item the API as the waiter
    \item the backend logic and database as the kitchen and storage
\end{itemize}

The customer does not walk into the kitchen and grab ingredients directly.
The customer makes a request through the waiter.
The waiter knows what requests are valid, forwards the work, then brings back the result.

That is exactly what your Ktor server is doing.

\subsection{What Ktor Is In That Story}

Ktor is the framework that helps you build the waiter:
\begin{itemize}
    \item it listens for requests
    \item it matches them to handlers
    \item it helps parse input
    \item it helps format output
    \item it lets you add cross-cutting behavior like auth, validation, logging, and rate limiting
\end{itemize}

\begin{kazecalloutbox}[How To Read The Rest]
If a chapter feels abstract, keep returning to this question: what is the client asking the server to do, and how does the server safely answer?
\end{kazecalloutbox}

\subsection{Important Technical Words You Will See Repeatedly}

Because this guide is for a beginner, here are some words that appear often in backend discussions.

\textbf{Server}

The program running on a machine that listens for requests and sends responses.

\textbf{Client}

The program that sends requests to the server. In your project, the mobile or web app is usually the client.

\textbf{Request}

A message sent by the client asking the server to do something.

\textbf{Response}

The answer sent back by the server.

\textbf{Endpoint}

A specific API address plus HTTP method, such as \texttt{GET /auth/me}.

\textbf{Route}

The Ktor code that handles an endpoint.

\textbf{JSON}

A text format used to exchange structured data between client and server.

\textbf{Plugin}

Reusable behavior added to the Ktor request/response pipeline, such as auth, JSON handling, or rate limiting.

\textbf{Middleware}

A more general web term for code that runs around request handling. In Ktor, the closest idea is usually a plugin.

\textbf{Dependency}

Anything a piece of code needs in order to do its job.

Examples:
\begin{itemize}
    \item a service depends on a repository
    \item a repository depends on a data source
    \item your project depends on Ktor or Koin libraries
\end{itemize}

\textbf{Dependency Injection}

A way of giving a class the things it needs from the outside instead of letting it build everything by itself.

\textbf{Coupling}

How tightly one piece of code is tied to another.

High coupling usually means code is harder to replace, test, or refactor.

\textbf{DTO}

Data Transfer Object. A class used to move data across a boundary such as HTTP.

\textbf{Business Logic}

The rules of the application itself, separate from HTTP details and database details.

\textbf{Header}

Extra metadata sent with a request or response, such as \texttt{Authorization} or \texttt{Content-Type}.

\textbf{Body}

The main content of a request or response. JSON is usually sent in the body.

\textbf{Status Code}

The numeric result of an HTTP response, such as:
\begin{itemize}
    \item \texttt{200 OK}
    \item \texttt{201 Created}
    \item \texttt{400 Bad Request}
    \item \texttt{401 Unauthorized}
    \item \texttt{404 Not Found}
    \item \texttt{500 Internal Server Error}
\end{itemize}

\textbf{Principal}

The authenticated identity attached to a request after auth succeeds.

\textbf{JWT Claim}

A piece of data stored inside a JWT, such as user id, email, or roles.

\textbf{Data Source}

An object that provides database connections to your application.

\textbf{Connection Pool}

A managed set of reusable database connections.

\textbf{Transaction}

A group of database operations that should succeed or fail together.

\textbf{Validation}

Checking whether incoming data is acceptable before deeper processing.

\textbf{Serialization}

Turning a Kotlin object into a transfer format such as JSON.

\textbf{Deserialization}

Turning JSON or another transfer format back into a Kotlin object.

\textbf{CORS}

Cross-Origin Resource Sharing. Browser rules that control which frontends are allowed to call your API.

\section{Course Project: A Tiny Library API You Can Imagine Building}

To make this document feel more like a course, we need one small project idea that keeps returning throughout the chapters.

Imagine we are building a tiny API for a library app.

The client app should be able to:
\begin{itemize}
    \item list books
    \item view one book
    \item add a new book
    \item sign a user in
    \item protect librarian-only actions
\end{itemize}

We will not fully implement the whole project in this document, but we will use it as the running teaching example.

\subsection{Why A Running Example Helps}

When you learn backend concepts in isolation, they can feel disconnected.

A mini project ties them together:
\begin{itemize}
    \item routing decides how \texttt{/books} requests are handled
    \item serialization decides how \texttt{BookDto} becomes JSON
    \item plugins handle validation, auth, and rate limiting
    \item repositories read and write books from storage
    \item tests verify the routes
    \item deployment makes the API available outside your machine
\end{itemize}

\subsection{The Minimal Shape Of The Example API}

\begin{longtable}{p{0.28\textwidth} p{0.64\textwidth}}
\textbf{Endpoint} & \textbf{Purpose} \\
\texttt{GET /books} & List books \\
\texttt{GET /books/\{id\}} & Get one book \\
\texttt{POST /books} & Create a book \\
\texttt{POST /auth/signin} & Sign a user in \\
\texttt{GET /auth/me} & Return the current signed-in user \\
\end{longtable}

\subsection{Why This Particular Example}

It is small enough for learning, but rich enough to include:
\begin{itemize}
    \item reading endpoints
    \item writing endpoints
    \item JSON
    \item validation
    \item database access
    \item authentication
    \item tests
\end{itemize}

That makes it a good beginner practice API.

\section{Chapter 1: Foundations \& Project Setup}

\subsection{What Ktor Is}

Ktor is a Kotlin framework for writing network applications. In this guide we focus on Ktor server, which means backend code that answers HTTP requests.

Ktor is not your database.
Ktor is not your business logic.
Ktor is not your app UI.

Ktor is the layer that receives requests, coordinates the pipeline, and returns responses.

\subsection{The Smallest Useful Server}

Here is a minimal Ktor application:

\begin{lstlisting}[style=kazekotlin]
fun main() {
    embeddedServer(Netty, port = 8080, module = Application::module)
        .start(wait = true)
}

fun Application.module() {
    routing {
        get("/") {
            call.respondText("Hello from Ktor")
        }
    }
}
\end{lstlisting}

Let us explain this like a course, line by line in concept:
\begin{itemize}
    \item \texttt{main()} is the JVM entry point
    \item \texttt{embeddedServer(...)} creates the web server
    \item \texttt{Netty} is the engine that actually handles network traffic
    \item \texttt{port = 8080} means the server listens on port 8080
    \item \texttt{Application::module} tells Ktor which function contains the app setup
    \item \texttt{start(wait = true)} starts the server and keeps it running
\end{itemize}

\subsection{What The Application Module Is}

The \texttt{module()} function is where you configure the app.

That usually includes:
\begin{itemize}
    \item installing plugins
    \item registering routes
    \item wiring dependencies
    \item connecting to the database
    \item loading configuration
\end{itemize}

So when you open a Ktor project and do not know where to begin, the application module is usually the best first stop.

\subsection{What Happens When A Request Arrives}

Suppose the browser or app requests \texttt{GET /}.

Ktor roughly does this:
\begin{enumerate}
    \item the server receives the request
    \item the application pipeline runs
    \item installed plugins may inspect or modify the request/response flow
    \item routing finds the matching route
    \item your handler code runs
    \item a response is sent back
\end{enumerate}

That means a Ktor app is not just ``route code.'' There is a full pipeline around the route.

\subsection{What To Pay Attention To As A Beginner}

\begin{itemize}
    \item Find the entry point first.
    \item Find the application module second.
    \item Separate setup code from business code.
    \item Notice what Ktor does for you automatically and what your own code is doing manually.
\end{itemize}

\subsection{If You Did Not Use A Framework}

Without Ktor, you would need to manually:
\begin{itemize}
    \item open sockets
    \item parse HTTP messages
    \item route paths
    \item parse headers and bodies
    \item construct HTTP responses
    \item manage pipeline behavior yourself
\end{itemize}

That is why frameworks exist.

\subsection{How Kaze Uses This}

Kaze's server startup file is \texttt{server/src/main/kotlin/dev/orestegabo/kaze/Application.kt}. It is more advanced than the tiny example, but it still follows the same idea: start the server, load configuration, install Koin, install Ktor features, then register routes.

\subsection{Advanced Kotlin Concepts Used Here}

\begin{itemize}
    \item \texttt{Application::module} is a function reference.
    \item \texttt{fun Application.module()} is an extension function.
    \item \texttt{?:} is the Elvis operator, used for fallback values.
\end{itemize}

\section{Chapter 2: Routing \& Request Handling}

\subsection{What Routing Is}

Routing is how the server decides which code should run for a request.

If a request comes in for:
\begin{itemize}
    \item \texttt{GET /books}
    \item \texttt{GET /books/42}
    \item \texttt{POST /books}
\end{itemize}

those are three different requests and may need three different handlers.

Routing is the map from request to code.

\subsection{A Minimal Example}

\begin{lstlisting}[style=kazekotlin]
fun Application.module() {
    routing {
        get("/books") {
            call.respondText("List all books")
        }

        get("/books/{id}") {
            val id = call.parameters["id"]
            call.respondText("Book id = $id")
        }

        post("/books") {
            call.respondText("Create a new book")
        }
    }
}
\end{lstlisting}

This teaches two important things:
\begin{itemize}
    \item the HTTP method matters
    \item the path matters
\end{itemize}

\texttt{GET /books} and \texttt{POST /books} are not the same action.

\subsection{What The HTTP Methods Usually Mean}

\begin{itemize}
    \item \texttt{GET}: read data
    \item \texttt{POST}: create something or trigger an action
    \item \texttt{PUT}: replace or update something
    \item \texttt{PATCH}: partially update something
    \item \texttt{DELETE}: delete something
\end{itemize}

These are conventions, but they are very important conventions. Good APIs follow them because it makes behavior predictable.

\subsection{What The \texttt{call} Object Is}

Inside a route, Ktor gives you \texttt{call}. Think of it as ``the current HTTP conversation.''

From \texttt{call}, you can access:
\begin{itemize}
    \item route parameters
    \item query parameters
    \item headers
    \item request body
    \item response helpers
\end{itemize}

Examples:

\begin{lstlisting}[style=kazekotlin]
val id = call.parameters["id"]
val page = call.request.queryParameters["page"]
val authHeader = call.request.headers["Authorization"]
\end{lstlisting}

\subsection{Path Parameters vs Query Parameters vs Request Body}

Beginners often mix these up, so let us make them very clear.

\textbf{Path parameters}

These identify a resource inside the URL path.

\begin{lstlisting}[style=kazekotlin]
get("/books/{id}") {
    val id = call.parameters["id"]
}
\end{lstlisting}

Example request:
\texttt{/books/42}

\textbf{Query parameters}

These modify or refine the request.

\begin{lstlisting}[style=kazekotlin]
get("/books") {
    val sort = call.request.queryParameters["sort"]
}
\end{lstlisting}

Example request:
\texttt{/books?sort=title}

\textbf{Request body}

This is the structured data sent with a request, usually JSON.

\begin{lstlisting}[style=kazekotlin]
@Serializable
data class CreateBookRequest(
    val title: String,
    val author: String
)

post("/books") {
    val request = call.receive<CreateBookRequest>()
}
\end{lstlisting}

Example JSON:

\begin{lstlisting}[style=kazekotlin]
{
  "title": "Effective Kotlin",
  "author": "Marcin Moskala"
}
\end{lstlisting}

\subsection{How A Good Route Should Feel}

A good route should usually be short.

Its job is:
\begin{enumerate}
    \item read input
    \item call a service
    \item respond
\end{enumerate}

Its job is usually not:
\begin{itemize}
    \item writing SQL directly
    \item containing complex business rules
    \item doing everything itself
\end{itemize}

\subsection{What Happens If Routes Do Too Much}

If routes contain too much logic:
\begin{itemize}
    \item they become hard to test
    \item they become hard to reuse
    \item business rules get mixed with HTTP details
    \item the code becomes fragile and harder to refactor
\end{itemize}

\subsection{Best Practices}

\begin{itemize}
    \item Keep routes thin.
    \item Use clear route naming.
    \item Prefer plural resource names for collections such as \texttt{/books}.
    \item Use the right HTTP method for the job.
    \item Separate public routes from protected ones.
\end{itemize}

\subsection{How Kaze Uses This}

Kaze groups routes under \texttt{/api/v1}, then further groups them by hotel, guest, and auth concerns. That is a good example of route structure following domain structure.

\subsection{Advanced Kotlin Concepts Used Here}

\begin{itemize}
    \item The routing API is a DSL using lambdas with receivers.
    \item \texttt{receive<Type>()} uses generics.
    \item Helper functions on \texttt{ApplicationCall} are often written as extension functions.
\end{itemize}

\section{Chapter 3: Content Negotiation \& Serialization}

\subsection{What Serialization Is}

Inside your Kotlin code, you work with Kotlin objects.

Across HTTP, clients and servers exchange text or bytes, usually JSON.

Serialization means:
\begin{quote}
turning a Kotlin object into JSON
\end{quote}

Deserialization means:
\begin{quote}
turning JSON into a Kotlin object
\end{quote}

\subsection{Why This Exists}

If your server returns a Kotlin data class directly, the client cannot magically understand the in-memory object.
The object has to be encoded into a format both sides agree on.

That is what JSON does.

\subsection{What Content Negotiation Means}

Content negotiation is the logic that helps decide how to encode and decode request and response bodies.

In simple terms:
\begin{itemize}
    \item the client says what kind of content it is sending or expects
    \item the server chooses the appropriate serializer/deserializer
\end{itemize}

In many Ktor apps, that mainly means JSON.

\subsection{Minimal Setup}

\begin{lstlisting}[style=kazekotlin]
install(ContentNegotiation) {
    json()
}
\end{lstlisting}

Once this is installed, you can do:

\begin{lstlisting}[style=kazekotlin]
@Serializable
data class BookDto(
    val id: String,
    val title: String
)

get("/books/1") {
    call.respond(BookDto("1", "Ktor Basics"))
}
\end{lstlisting}

Ktor will serialize the \texttt{BookDto} to JSON for you.

\subsection{Why \texttt{@Serializable} Is Needed}

\texttt{@Serializable} tells Kotlin's serialization library that the class should be encodable and decodable.

Without it, Ktor usually will not know how to serialize or deserialize the class properly.

\subsection{Receiving JSON}

\begin{lstlisting}[style=kazekotlin]
@Serializable
data class CreateBookRequest(
    val title: String,
    val author: String
)

post("/books") {
    val request = call.receive<CreateBookRequest>()
    call.respondText("Created ${request.title}")
}
\end{lstlisting}

What happens here:
\begin{enumerate}
    \item the client sends JSON
    \item Ktor reads the body
    \item Ktor deserializes the JSON into \texttt{CreateBookRequest}
    \item your route receives a normal Kotlin object
\end{enumerate}

\subsection{What Happens If You Do Not Install Content Negotiation}

Without this plugin:
\begin{itemize}
    \item you may need to manually read request bodies
    \item you may need to manually build JSON strings
    \item your code becomes more repetitive and error-prone
\end{itemize}

\subsection{DTOs and Why They Matter}

DTO means Data Transfer Object.

These are classes made specifically for moving data across boundaries such as HTTP.

Why use DTOs:
\begin{itemize}
    \item not every internal field should be exposed
    \item request shape and domain shape are often different
    \item API contracts should be explicit
\end{itemize}

Best practice:
\begin{itemize}
    \item use request DTOs for incoming JSON
    \item use response DTOs for outgoing JSON
    \item keep domain models separate when needed
\end{itemize}

\subsection{How Kaze Uses This}

Kaze installs \texttt{ContentNegotiation} in \texttt{ApiSupport.kt} and uses DTO-based request and response handling across auth and business routes.

\subsection{Advanced Kotlin Concepts Used Here}

\begin{itemize}
    \item Data classes are especially good for DTOs.
    \item Builders like \texttt{Json \{ ... \}} are Kotlin DSLs.
    \item Generic APIs like \texttt{receive<Type>()} are a common Kotlin technique.
\end{itemize}

\section{Chapter 4: Plugins, Middleware \& Rate Limiting}

\subsection{What A Plugin Is}

This chapter is one of the most important because you specifically asked about plugins, middleware, and rate limiting.

In Ktor, a plugin is reusable behavior that plugs into the request/response pipeline.

Plain-language version:
\begin{quote}
A plugin is code that runs around your routes to handle common concerns.
\end{quote}

Examples of common concerns:
\begin{itemize}
    \item logging
    \item authentication
    \item JSON handling
    \item validation
    \item error formatting
    \item response compression
    \item rate limiting
\end{itemize}

\subsection{What Middleware Means}

Different ecosystems use different words.

In many web frameworks, people say \textbf{middleware}.
In Ktor, the official word is usually \textbf{plugin}.

The ideas are very close:
\begin{itemize}
    \item code runs before the route
    \item code may run after the route
    \item code may inspect, reject, modify, or enrich the request/response process
\end{itemize}

So if you hear someone say ``middleware'' while discussing Ktor, they usually mean plugin-like pipeline behavior.

\subsection{How To Mentally Picture Plugins}

Imagine a request entering the server like a package moving through checkpoints:
\begin{enumerate}
    \item logging checkpoint
    \item CORS checkpoint
    \item authentication checkpoint
    \item validation checkpoint
    \item route handler
    \item error handling / response formatting checkpoint
\end{enumerate}

That is not the exact internal Ktor implementation, but it is a very good beginner mental model.

\subsection{Why Plugins Exist}

Without plugins, every route would need to repeat the same work:
\begin{itemize}
    \item log the request
    \item catch exceptions
    \item check auth
    \item validate input
    \item compress output
\end{itemize}

That would make the code repetitive and inconsistent.

Plugins centralize those repeated concerns.

\subsection{Common Plugins And What They Do}

\textbf{ContentNegotiation}

Handles JSON serialization/deserialization.

\textbf{CallLogging}

Logs requests so you can see what the server is doing.

\textbf{StatusPages}

Transforms exceptions into clean HTTP responses.

\textbf{Authentication}

Protects routes and identifies the caller.

\textbf{RequestValidation}

Rejects bad inputs early.

\textbf{Compression}

Compresses responses to save bandwidth.

\textbf{CORS}

Controls which browser origins are allowed to call the API.

\textbf{RateLimit}

Restricts how many requests a client can make in a time window.

\subsection{What Rate Limiting Is}

Rate limiting is a protective rule like:
\begin{quote}
this client can only make 100 requests per minute
\end{quote}

Why this matters:
\begin{itemize}
    \item stops accidental spam from buggy clients
    \item slows abuse
    \item helps protect infrastructure
    \item is useful on login and sensitive endpoints
\end{itemize}

\subsection{How Rate Limiting Works Conceptually}

Any rate limiter needs at least two decisions:
\begin{itemize}
    \item \textbf{who} are we counting requests for
    \item \textbf{how many} requests are allowed in what period
\end{itemize}

The ``who'' is usually a key such as:
\begin{itemize}
    \item IP address
    \item user id
    \item API key
    \item session id
\end{itemize}

The ``how many'' is usually a limit and a window such as:
\begin{itemize}
    \item 60 requests per minute
    \item 1000 requests per hour
\end{itemize}

\subsection{What Happens If You Do Not Use Rate Limiting}

Without rate limiting:
\begin{itemize}
    \item a broken client may hammer your API
    \item brute-force login attempts are easier
    \item expensive endpoints can be abused
    \item one client can consume unfair resources
\end{itemize}

\subsection{What To Pay Attention To}

\begin{itemize}
    \item Choose the right key. IP-based limits are simple but imperfect.
    \item Choose limits that protect the app without hurting normal users.
    \item Apply stricter limits to sensitive routes like auth.
    \item Remember that rate limiting is protection, not a substitute for proper auth or validation.
\end{itemize}

\subsection{Best Practices For Plugins}

\begin{itemize}
    \item Put cross-cutting concerns in plugins, not inside every route.
    \item Keep plugin setup centralized when possible.
    \item Validate inputs early.
    \item Translate exceptions centrally.
    \item Do not over-install things you do not understand; know why each plugin exists.
\end{itemize}

\subsection{How Kaze Uses This}

Kaze installs many plugins in \texttt{ApiSupport.kt}: content negotiation, authentication, request validation, rate limiting, compression, and status pages. Once you know the concepts, that file becomes much easier to read because you can say ``this is the shared pipeline behavior'' instead of seeing it as random setup code.

\subsection{Advanced Kotlin Concepts Used Here}

\begin{itemize}
    \item Plugin configuration uses lambdas and builder DSLs.
    \item Small utilities like \texttt{takeIf} and safe calls are common in config code.
    \item Centralized config objects are often immutable data classes.
\end{itemize}

\begin{kazehighlightbox}
\textbf{Learning goal for this chapter:} a route is not the whole story. Plugins are the shared behavior around routes, and rate limiting is one example of defensive shared behavior.
\end{kazehighlightbox}

\section{Chapter 5: Advanced Database Integration}

\subsection{What The Backend Needs From Storage}

Most useful APIs need persistence. That means the server must be able to store and retrieve data.

Typical needs:
\begin{itemize}
    \item users
    \item sessions
    \item business objects such as hotels, events, and bookings
    \item logs or audit information
\end{itemize}

\subsection{Why Ktor Does Not ``Do The Database'' For You}

Ktor is a web framework, not a database framework.

That means you combine it with other libraries and patterns for storage:
\begin{itemize}
    \item JDBC
    \item Exposed
    \item JPA/Hibernate
    \item jOOQ
    \item custom repositories
\end{itemize}

\subsection{The Repository Pattern}

Think of a repository as the storage translator for the rest of your app.

Routes should not know SQL.
Services should ideally not care about low-level JDBC details.

Example interface:

\begin{lstlisting}[style=kazekotlin]
interface BookRepository {
    suspend fun listBooks(): List<Book>
    suspend fun findBook(id: String): Book?
    suspend fun saveBook(book: Book): Book
}
\end{lstlisting}

\subsection{Why Connection Pools Matter}

Databases are expensive resources.

You do not want to open a brand new database connection from scratch for every request and throw it away carelessly.

A connection pool:
\begin{itemize}
    \item keeps reusable connections available
    \item improves performance
    \item avoids too many simultaneous connection creations
    \item gives centralized timeout and pool settings
\end{itemize}

Kaze uses HikariCP for this.

\subsection{What JDBC Looks Like}

If you have not worked with backend storage much, raw JDBC can look noisy. That is normal.

Basic JDBC flow:
\begin{enumerate}
    \item get a connection
    \item prepare a statement
    \item bind parameters
    \item execute query or update
    \item read results
    \item close resources
\end{enumerate}

Example:

\begin{lstlisting}[style=kazekotlin]
dataSource.connection.use { connection ->
    connection.prepareStatement(
        "SELECT id, title FROM books WHERE id = ?"
    ).use { statement ->
        statement.setString(1, bookId)
        statement.executeQuery().use { result ->
            if (result.next()) {
                Book(
                    id = result.getString("id"),
                    title = result.getString("title")
                )
            } else null
        }
    }
}
\end{lstlisting}

\subsection{Why Prepared Statements Matter}

Prepared statements are important because they:
\begin{itemize}
    \item reduce SQL injection risk
    \item separate SQL structure from parameter values
    \item are the standard safe way to pass user input into queries
\end{itemize}

\subsection{What Happens If You Mix SQL Into Routes}

If routes directly contain database code:
\begin{itemize}
    \item your HTTP layer becomes tightly coupled to storage
    \item testing becomes harder
    \item reuse gets worse
    \item the code becomes harder to organize
\end{itemize}

\subsection{Best Practices}

\begin{itemize}
    \item Keep SQL or storage details in repositories.
    \item Use prepared statements.
    \item Use connection pooling.
    \item Think carefully about transactions for multi-step writes.
    \item Keep database setup and migration strategy explicit.
\end{itemize}

\subsection{How Kaze Uses This}

Kaze loads database config, creates a Hikari-backed data source, initializes schema and seed data, and uses JDBC-heavy repositories to access PostgreSQL data.

\subsection{Advanced Kotlin Concepts Used Here}

\begin{itemize}
    \item \texttt{use \{ ... \}} handles resource cleanup.
    \item \texttt{buildList \{ ... \}} is useful for building immutable result lists.
    \item \texttt{suspend} on repository methods does not automatically make blocking JDBC non-blocking.
\end{itemize}

\section{Chapter 6: Authentication \& JWT Security (Current)}

\subsection{What Authentication Is}

Authentication means proving identity.

Examples:
\begin{itemize}
    \item email and password
    \item Google sign-in
    \item Apple sign-in
    \item API key
\end{itemize}

The main question is:
\begin{quote}
who is making this request?
\end{quote}

\subsection{What Authorization Is}

Authorization is different:
\begin{quote}
now that we know who you are, what are you allowed to do?
\end{quote}

Example:
\begin{itemize}
    \item any signed-in user can read their own profile
    \item only admins can view admin routes
    \item only invited users can access some event data
\end{itemize}

\subsection{What A JWT Is}

JWT means JSON Web Token.

It is a signed token that can carry claims, such as:
\begin{itemize}
    \item user id
    \item email
    \item roles
    \item expiration time
\end{itemize}

The key idea:
\begin{itemize}
    \item the server signs the token
    \item later, the server can verify that the token was not tampered with
\end{itemize}

\subsection{Why JWTs Are Common In Mobile/API Systems}

JWTs are popular because:
\begin{itemize}
    \item clients can store them and send them as bearer tokens
    \item the server can verify them efficiently
    \item they fit well with HTTP APIs
\end{itemize}

But they are not automatically secure just because they are JWTs. Design still matters.

\subsection{Access Tokens And Refresh Tokens}

Best practice is often:
\begin{itemize}
    \item short-lived access token
    \item longer-lived refresh token
\end{itemize}

Why:
\begin{itemize}
    \item if an access token is stolen, its lifetime is limited
    \item refresh tokens let the client renew the session without forcing frequent logins
\end{itemize}

\subsection{How Ktor Protects Routes}

Ktor uses the \texttt{Authentication} plugin.

Basic flow:
\begin{enumerate}
    \item install the auth mechanism
    \item wrap protected routes in \texttt{authenticate \{ ... \}}
    \item read the principal inside the route
\end{enumerate}

Conceptual example:

\begin{lstlisting}[style=kazekotlin]
authenticate("auth-jwt") {
    get("/me") {
        val principal = call.principal<JWTPrincipal>()
        call.respondText("User = ${principal?.payload?.subject}")
    }
}
\end{lstlisting}

\subsection{What Happens If You Do Not Protect Sensitive Routes}

If you forget auth or authorization:
\begin{itemize}
    \item private data may leak
    \item users may modify things they should not
    \item your API becomes unsafe even if the code ``works''
\end{itemize}

\subsection{Password Hashing}

Never store plain passwords.

Use password hashing such as bcrypt. Kaze uses bcrypt. That means:
\begin{itemize}
    \item the server stores a hash, not the original password
    \item login checks the password against the hash
\end{itemize}

\subsection{Best Practices}

\begin{itemize}
    \item Keep access tokens short-lived.
    \item Store refresh tokens carefully.
    \item Hash passwords.
    \item Separate authentication from authorization in your thinking.
    \item Protect only the routes that need protection, but do not forget sensitive ones.
\end{itemize}

\subsection{How Kaze Uses This}

Kaze currently supports password auth, social sign-in, JWT access tokens, refresh tokens, and protected routes such as \texttt{/auth/me}. That makes it a strong real example after you understand the principles.

\subsection{Advanced Kotlin Concepts Used Here}

\begin{itemize}
    \item Kotlin null-safe chains are common in auth normalization.
    \item \texttt{when} is often used as a clean branching expression.
    \item Custom properties with getters can lazily resolve dependencies.
\end{itemize}

\section{Chapter 7: Dependency Injection With Koin}

\subsection{What A Dependency Is}

This is one of the most important words in backend development, so it is worth defining very clearly.

A dependency is anything a piece of code needs in order to work.

Examples:
\begin{itemize}
    \item a service may need a repository
    \item a repository may need a database connection or data source
    \item an auth service may need JWT configuration
    \item your whole project may need external libraries such as Ktor, Koin, or the PostgreSQL driver
\end{itemize}

So if a class cannot do its work alone and needs another object, config, or library, that thing is a dependency.

\subsection{Two Meanings Of Dependency}

You will hear the word ``dependency'' in two slightly different ways.

\textbf{1. Code dependency}

This means one class or component needs another class or component.

Example:

\begin{lstlisting}[style=kazekotlin]
class BookService(
    private val repository: BookRepository
)
\end{lstlisting}

Here, \texttt{BookService} depends on \texttt{BookRepository}.

\textbf{2. Project or library dependency}

This means your project needs an external library in order to compile or run.

Examples in Gradle:

\begin{lstlisting}[style=kazekotlin]
implementation(libs.ktor.serverCore)
implementation(libs.koin.ktor)
implementation(libs.postgresql)
\end{lstlisting}

These are library dependencies.

\subsection{Why Dependencies Are Needed}

Real applications are built from smaller pieces with different responsibilities.

For example:
\begin{itemize}
    \item routes handle HTTP
    \item services handle business logic
    \item repositories handle storage
    \item config objects hold configuration
\end{itemize}

Because one layer needs help from another layer, dependencies naturally appear.

That is not a problem. In fact, it is normal and healthy.

The real question is not:
\begin{quote}
can we avoid dependencies completely?
\end{quote}

The real question is:
\begin{quote}
are dependencies organized clearly and safely?
\end{quote}

\subsection{What Goes Wrong When Dependencies Are Poorly Managed}

If a class creates everything by itself, several problems appear:
\begin{itemize}
    \item the class does too much
    \item the class becomes harder to test
    \item replacing one part becomes difficult
    \item code gets tightly coupled
\end{itemize}

Example of a class doing too much:

\begin{lstlisting}[style=kazekotlin]
class BookService {
    fun listBooks(): List<Book> {
        val dataSource = createDataSource()
        val repository = JdbcBookRepository(dataSource)
        return repository.listBooks()
    }
}
\end{lstlisting}

This works, but it mixes object creation with business behavior.

Better:

\begin{lstlisting}[style=kazekotlin]
class BookService(
    private val repository: BookRepository
) {
    fun listBooks(): List<Book> = repository.listBooks()
}
\end{lstlisting}

Now the service focuses on its real job.

\subsection{What Coupling Means}

Coupling describes how tightly one piece of code is tied to another.

If a class is tightly coupled to a specific implementation, it becomes harder to:
\begin{itemize}
    \item replace that implementation
    \item test the class in isolation
    \item reuse the class elsewhere
\end{itemize}

So when we talk about dependency injection, one big goal is reducing unnecessary coupling.

\subsection{What Dependency Injection Really Means}

Dependency injection sounds fancy, but the core idea is simple:

\begin{quote}
objects should receive the things they need instead of creating everything by themselves
\end{quote}

If a service needs a repository, it should usually be given that repository rather than constructing it internally.

\subsection{Why This Matters}

Without DI:
\begin{itemize}
    \item classes become tightly coupled
    \item object creation gets repeated everywhere
    \item testing gets harder
    \item wiring gets hidden inside business code
\end{itemize}

With DI:
\begin{itemize}
    \item wiring is centralized
    \item classes state their dependencies clearly
    \item test doubles become easier
    \item architecture becomes easier to reason about
\end{itemize}

\subsection{Constructor Injection}

This is the clean beginner-friendly pattern:

\begin{lstlisting}[style=kazekotlin]
class BookService(
    private val repository: BookRepository
)
\end{lstlisting}

The service says plainly what it depends on.

\subsection{What Koin Does}

Koin is a DI framework for Kotlin. It helps you register objects and ask for them later.

Conceptual example:

\begin{lstlisting}[style=kazekotlin]
val appModule = module {
    single<BookRepository> { JdbcBookRepository(get()) }
    single { BookService(get()) }
}
\end{lstlisting}

Meaning:
\begin{itemize}
    \item create one shared \texttt{BookRepository}
    \item create one shared \texttt{BookService}
    \item if a constructor needs another object, resolve it with \texttt{get()}
\end{itemize}

\subsection{What Happens If You Do Not Use DI}

You do not absolutely need a DI framework for every project.

But as projects grow, without some DI pattern you often get:
\begin{itemize}
    \item scattered setup logic
    \item repeated manual object construction
    \item harder testing
    \item hidden dependencies
\end{itemize}

\subsection{Best Practices}

\begin{itemize}
    \item Prefer constructor injection.
    \item Keep wiring centralized.
    \item Depend on interfaces where useful.
    \item Do not use DI to hide bad design; use it to clarify assembly.
    \item Separate object creation from business behavior.
\end{itemize}

\subsection{How Kaze Uses This}

Kaze installs Koin in the server module and registers config, data source, repositories, auth helpers, services, and grouped server dependencies.

\subsection{Advanced Kotlin Concepts Used Here}

\begin{itemize}
    \item Koin uses DSL-style builders.
    \item Generic bindings let interfaces point to concrete implementations.
    \item Constructor injection reads especially cleanly in Kotlin.
\end{itemize}

\section{Chapter 8: Testing \& Error Handling}

\subsection{Why Testing A Backend Is Different}

In backend work, you are not only testing functions. You are also testing:
\begin{itemize}
    \item HTTP behavior
    \item route registration
    \item auth rules
    \item validation
    \item status codes
    \item integration with storage or configuration
\end{itemize}

\subsection{What Ktor Testing Gives You}

Ktor lets you start the application in a test environment and send requests to it.

Conceptually:

\begin{lstlisting}[style=kazekotlin]
testApplication {
    application { module() }

    val response = client.get("/health")
    assertEquals(HttpStatusCode.OK, response.status)
}
\end{lstlisting}

This is extremely useful because it tests the app close to how it really behaves.

\subsection{What Error Handling Means On A Server}

Good server error handling means:
\begin{itemize}
    \item bad client input gets a clear client error
    \item unauthorized access gets the proper auth error
    \item missing resources get not-found responses
    \item unexpected failures get safe server errors
\end{itemize}

This is important because APIs are contracts. They should fail in understandable ways, not with random stack traces or vague broken behavior.

\subsection{Status Pages}

Ktor's \texttt{StatusPages} plugin is the central place for translating exceptions into responses.

Why that is good:
\begin{itemize}
    \item routes stay cleaner
    \item error responses stay consistent
    \item you avoid repetitive try/catch blocks everywhere
\end{itemize}

\subsection{What Happens If You Ignore Error Handling}

Without deliberate error handling:
\begin{itemize}
    \item clients get inconsistent responses
    \item debugging gets harder
    \item sensitive internal details might leak
    \item the API becomes unpleasant to use
\end{itemize}

\subsection{Best Practices}

\begin{itemize}
    \item Test happy paths and failure paths.
    \item Test authentication and validation behavior, not only business success cases.
    \item Keep error responses consistent.
    \item Use centralized exception translation.
\end{itemize}

\subsection{How Kaze Uses This}

Kaze tests health checks, Swagger, auth rules, validation failures, rate limiting, and compression. It also uses \texttt{StatusPages} to convert exceptions into structured API error responses.

\subsection{Advanced Kotlin Concepts Used Here}

\begin{itemize}
    \item Test builders often use suspending lambdas with receivers.
    \item \texttt{vararg} is useful for flexible test config helpers.
    \item Companion objects keep local test constants tidy.
\end{itemize}

\section{Chapter 9: Deployment, Docker, and KMP Integration}

\subsection{Why Deployment Is Part Of Learning Backend}

If your code only runs in your IDE, you have learned only part of backend development.

A real backend must also be:
\begin{itemize}
    \item built
    \item configured
    \item packaged
    \item deployed
    \item run in another environment
\end{itemize}

\subsection{What Docker Solves}

Docker packages the application with its runtime environment into a container image.

Why people use it:
\begin{itemize}
    \item consistent runtime environment
    \item easier deployment
    \item cleaner separation between build stage and runtime stage
    \item fewer machine-specific surprises
\end{itemize}

Kaze uses a multi-stage Docker build, which is a strong production pattern.

\subsection{What Configuration Means In Deployment}

A backend needs runtime configuration such as:
\begin{itemize}
    \item database URL
    \item credentials
    \item JWT secret
    \item OAuth secrets
    \item host or CORS settings
\end{itemize}

That is why environment variables matter so much.

\subsection{What KMP Changes}

Kotlin Multiplatform means your project can share some Kotlin code across platforms.

But that does \emph{not} mean the backend and app are the same thing.

In Kaze:
\begin{itemize}
    \item the server is a Ktor JVM module
    \item the shared module contains common logic and models
    \item the Compose app is the client
\end{itemize}

\subsection{What To Pay Attention To}

\begin{itemize}
    \item shared language does not remove the network boundary
    \item clients and servers still communicate through contracts
    \item deployment of the server is separate from shipping the mobile app
\end{itemize}

\subsection{Best Practices}

\begin{itemize}
    \item keep secrets out of source code
    \item make config explicit
    \item use stable deployment scripts
    \item keep the backend independently runnable
\end{itemize}

\subsection{How Kaze Uses This}

Kaze packages the backend in Docker and deploys it with a Cloud Run script that passes environment configuration. On the client side, the Compose app uses a Ktor HTTP client for auth calls. That is a clean example of a real KMP app talking to a real Ktor backend.

\subsection{Advanced Kotlin Concepts Used Here}

\begin{itemize}
    \item KMP uses \texttt{expect}/\texttt{actual} for platform-specific implementations.
    \item Mapping helpers convert transport-layer objects into app-layer models.
    \item Immutable config objects are common in deployment-aware code.
\end{itemize}

\section{Final Advice For Studying Kaze}

Do not start by reading the whole Kaze backend end-to-end.

Instead:
\begin{enumerate}
    \item understand what an API is
    \item understand routes
    \item understand JSON request/response handling
    \item understand plugins
    \item understand auth and storage separately
    \item only then read the bigger project files
\end{enumerate}

Best reading order in the project after this guide:
\begin{enumerate}
    \item \texttt{Application.kt}
    \item \texttt{ApiSupport.kt}
    \item \texttt{ApiRouting.kt}
    \item \texttt{AuthRouting.kt}
    \item \texttt{AuthService.kt}
    \item \texttt{ServerModules.kt}
    \item \texttt{DatabaseConfig.kt}
    \item \texttt{JdbcRepositories.kt}
    \item \texttt{ApiRoutesTest.kt}
\end{enumerate}

\begin{kazehighlightbox}
\textbf{Most important thing to remember:} Ktor is not hard because of syntax. It feels hard when too many backend concepts arrive at the same time. Once you separate the ideas --- routing, JSON, plugins, auth, database, DI, testing, deployment --- the code becomes far more understandable.
\end{kazehighlightbox}

\section{Visual Mental Models}

This section is here to make the concepts easier to picture.

\subsection{The Basic Request Flow}

\begin{longtable}{p{0.2\textwidth} p{0.72\textwidth}}
\textbf{Step 1} & Client sends an HTTP request such as \texttt{GET /books}. \\
\textbf{Step 2} & Ktor server receives the request. \\
\textbf{Step 3} & Plugins may inspect or modify the request pipeline. \\
\textbf{Step 4} & Routing matches the request to a route handler. \\
\textbf{Step 5} & The route calls services and repositories as needed. \\
\textbf{Step 6} & The route creates or receives a result. \\
\textbf{Step 7} & Ktor serializes the response and sends it back. \\
\end{longtable}

Short version:

\begin{quote}
Client -> Ktor pipeline -> Route -> Service -> Repository -> Database -> back to client
\end{quote}

\subsection{The Plugin Mental Model}

\begin{longtable}{p{0.2\textwidth} p{0.72\textwidth}}
\textbf{Before Route} & Logging, CORS, authentication checks, request validation may happen before your route code runs. \\
\textbf{Route} & Your route reads input and coordinates the main action. \\
\textbf{After Route} & Serialization, compression, response headers, and centralized error formatting may affect the outgoing response. \\
\end{longtable}

\subsection{The Login Flow Mental Model}

\begin{longtable}{p{0.2\textwidth} p{0.72\textwidth}}
\textbf{Step 1} & Client sends credentials to a sign-in endpoint. \\
\textbf{Step 2} & Server verifies identity. \\
\textbf{Step 3} & Server creates a token or session. \\
\textbf{Step 4} & Client stores the access token. \\
\textbf{Step 5} & Client sends the token on future protected requests. \\
\textbf{Step 6} & Server validates the token before allowing the route. \\
\end{longtable}

\subsection{The Refresh Token Mental Model}

\begin{longtable}{p{0.2\textwidth} p{0.72\textwidth}}
\textbf{Step 1} & Access token expires. \\
\textbf{Step 2} & Client sends refresh token to a refresh endpoint. \\
\textbf{Step 3} & Server checks whether the refresh token is still valid. \\
\textbf{Step 4} & Server issues a new access token, and often a replacement refresh token. \\
\textbf{Step 5} & Client continues using the new access token. \\
\end{longtable}

\section{Common Mistakes, Symptoms, and Debugging Hints}

Beginners usually learn fastest when they can connect mistakes to symptoms.

\subsection{Routing Mistakes}

\begin{longtable}{p{0.26\textwidth} p{0.28\textwidth} p{0.36\textwidth}}
\textbf{Mistake} & \textbf{Symptom} & \textbf{What To Check} \\
Wrong HTTP method & \texttt{404} or route not reached & Check whether the route is \texttt{GET}, \texttt{POST}, etc. \\
Wrong path & \texttt{404} & Check exact path spelling and prefixes like \texttt{/api/v1}. \\
Missing path parameter & \texttt{400} or crash & Check route template and parameter extraction. \\
\end{longtable}

\subsection{Serialization Mistakes}

\begin{longtable}{p{0.26\textwidth} p{0.28\textwidth} p{0.36\textwidth}}
\textbf{Mistake} & \textbf{Symptom} & \textbf{What To Check} \\
Forgot \texttt{ContentNegotiation} & JSON handling fails or feels manual & Verify the plugin is installed. \\
Forgot \texttt{@Serializable} & Serialization or deserialization errors & Add the annotation to DTO classes. \\
Wrong request JSON shape & \texttt{400} or parsing failure & Compare client JSON fields with DTO fields. \\
\end{longtable}

\subsection{Plugin and Auth Mistakes}

\begin{longtable}{p{0.26\textwidth} p{0.28\textwidth} p{0.36\textwidth}}
\textbf{Mistake} & \textbf{Symptom} & \textbf{What To Check} \\
Protected route not wrapped in \texttt{authenticate} & Route is accidentally public & Check route grouping. \\
Bad or missing token & \texttt{401 Unauthorized} & Check header name, token value, and verifier config. \\
Validation missing & Bad data reaches services & Add \texttt{RequestValidation} or route-level checks. \\
\end{longtable}

\subsection{Database Mistakes}

\begin{longtable}{p{0.26\textwidth} p{0.28\textwidth} p{0.36\textwidth}}
\textbf{Mistake} & \textbf{Symptom} & \textbf{What To Check} \\
Wrong DB config & Startup or query failures & Check URL, user, password, and environment variables. \\
Mixing SQL into routes & Messy route code & Move storage logic into repositories. \\
Unsafe SQL building & Security risk or broken queries & Use prepared statements. \\
\end{longtable}

\subsection{A Debugging Habit That Helps}

When something breaks, ask these questions in order:
\begin{enumerate}
    \item Did the request reach the correct route?
    \item Was the request data parsed correctly?
    \item Did a plugin reject it first?
    \item Did the service logic fail?
    \item Did the repository or database fail?
    \item Was the response serialized correctly?
\end{enumerate}

\section{Exercises}

These are intentionally small. The goal is to make you think like a backend developer without overwhelming you.

\subsection{After Chapter 1}

\begin{itemize}
    \item Write the smallest Ktor server you can that returns \texttt{"hello"} on \texttt{GET /}.
    \item Change the port from \texttt{8080} to another number and explain what changed.
\end{itemize}

\subsection{After Chapter 2}

\begin{itemize}
    \item Design three routes for a movie API.
    \item Add one path parameter route and one query parameter route.
    \item Explain in your own words when to use request body data instead of query parameters.
\end{itemize}

\subsection{After Chapter 3}

\begin{itemize}
    \item Create a \texttt{BookDto} class with \texttt{@Serializable}.
    \item Write a route that returns one \texttt{BookDto}.
    \item Write a request DTO for creating a book.
\end{itemize}

\subsection{After Chapter 4}

\begin{itemize}
    \item List three concerns that belong in plugins instead of route code.
    \item Explain rate limiting in one sentence and give one route that should probably be rate-limited.
    \item Explain what would happen if a backend had no validation and no status-page style error handling.
\end{itemize}

\subsection{After Chapter 5}

\begin{itemize}
    \item Write a repository interface for books or tasks.
    \item Explain why prepared statements are safer than SQL string concatenation.
    \item Explain what a connection pool is in your own words.
\end{itemize}

\subsection{After Chapter 6}

\begin{itemize}
    \item Explain the difference between authentication and authorization.
    \item Explain why access tokens are usually short-lived.
    \item Explain why passwords should be hashed, not stored directly.
\end{itemize}

\subsection{After Chapter 7}

\begin{itemize}
    \item Identify one code dependency and one library dependency in your own words.
    \item Rewrite a class constructor to receive a dependency instead of creating it internally.
    \item Explain coupling in one short paragraph.
\end{itemize}

\subsection{After Chapter 8}

\begin{itemize}
    \item Write a simple test idea for a \texttt{GET /health} endpoint.
    \item List two failure cases you would test for an auth route.
\end{itemize}

\subsection{After Chapter 9}

\begin{itemize}
    \item Explain why environment variables are useful in deployment.
    \item Explain why a mobile app and backend are still separate systems even if both are written in Kotlin.
\end{itemize}

\section{Java and Spring Bridges}

Because you already know Java, it helps to connect Ktor ideas to a more traditional Java backend mindset.

\begin{longtable}{p{0.28\textwidth} p{0.28\textwidth} p{0.34\textwidth}}
\textbf{Ktor Idea} & \textbf{Rough Java/Spring Analogy} & \textbf{Important Difference} \\
Route & Controller endpoint & Ktor often feels lighter and more DSL-driven. \\
Plugin & Filter / interceptor / middleware-like behavior & Ktor uses pipeline/plugin vocabulary. \\
Application module & Application setup / configuration class & Ktor often centralizes setup in functions instead of heavy annotations. \\
Koin module & Spring bean configuration & Koin is lighter and more explicit in code. \\
DTO & Request/response DTO & Same basic idea. \\
Repository & Repository / DAO & Same basic architectural role. \\
\end{longtable}

\subsection{One Important Mindset Difference}

Ktor often feels more explicit and less annotation-heavy than frameworks many Java developers first meet.

That can feel unusual at first, but it is often helpful for learning because more of the setup is visible in code.

\section{Kotlin for Backend Appendix}

This appendix exists because backend Kotlin uses some features that may not be familiar if you mainly coded Java or C++.

\subsection{Extension Functions}

Example:

\begin{lstlisting}[style=kazekotlin]
fun Application.module() { ... }
\end{lstlisting}

This adds a function to an existing type without subclassing it.

\subsection{Data Classes}

These are concise classes often used for DTOs and config objects.

\begin{lstlisting}[style=kazekotlin]
data class BookDto(
    val id: String,
    val title: String
)
\end{lstlisting}

\subsection{Null Safety}

Kotlin makes null handling explicit with tools like:
\begin{itemize}
    \item safe call: \texttt{?.}
    \item Elvis operator: \texttt{?:}
    \item nullable types such as \texttt{String?}
\end{itemize}

\subsection{Lambdas With Receivers}

Ktor and Koin DSLs use this heavily. That is why blocks like:

\begin{lstlisting}[style=kazekotlin]
routing {
    get("/") { ... }
}
\end{lstlisting}

read like mini languages.

\subsection{Generics}

You will see this in APIs like:

\begin{lstlisting}[style=kazekotlin]
call.receive<CreateBookRequest>()
\end{lstlisting}

The function works with different types.

\subsection{\texttt{suspend}}

\texttt{suspend} marks functions that can be called from coroutines. It is a key part of modern Kotlin backend programming, but remember: \texttt{suspend} does not automatically mean ``non-blocking'' for everything inside the function.

\subsection{\texttt{object} and Companion Object}

\texttt{object} creates a singleton-like declaration.

\texttt{companion object} is often used for constants or factory-like helpers associated with a class.

\section{Mapping The Concepts Back To Kaze}

Now that the course material has taught the ideas separately, here is how they map back to your project.

\begin{longtable}{p{0.28\textwidth} p{0.64\textwidth}}
\textbf{Concept} & \textbf{Where To Look In Kaze} \\
Server startup and application module & \texttt{server/src/main/kotlin/dev/orestegabo/kaze/Application.kt} \\
Ktor plugin setup & \texttt{server/src/main/kotlin/dev/orestegabo/kaze/api/ApiSupport.kt} \\
Main route structure & \texttt{server/src/main/kotlin/dev/orestegabo/kaze/api/ApiRouting.kt} \\
Auth-specific routes & \texttt{server/src/main/kotlin/dev/orestegabo/kaze/api/AuthRouting.kt} \\
Auth logic & \texttt{server/src/main/kotlin/dev/orestegabo/kaze/auth/AuthService.kt} \\
Dependency injection & \texttt{server/src/main/kotlin/dev/orestegabo/kaze/di/ServerModules.kt} \\
Database config and connection setup & \texttt{server/src/main/kotlin/dev/orestegabo/kaze/infrastructure/DatabaseConfig.kt} and \texttt{DatabaseFactory.kt} \\
Repository-style storage access & \texttt{server/src/main/kotlin/dev/orestegabo/kaze/infrastructure/JdbcRepositories.kt} \\
Endpoint tests & \texttt{server/src/test/kotlin/dev/orestegabo/kaze/ApiRoutesTest.kt} \\
Client talking to backend & \texttt{composeApp/src/commonMain/kotlin/dev/orestegabo/kaze/presentation/auth/AuthGateway.kt} \\
\end{longtable}

\subsection{Best Way To Re-read Kaze After This}

Once the course concepts feel familiar, try re-reading the project with this sequence:
\begin{enumerate}
    \item startup
    \item plugins
    \item routes
    \item auth
    \item DI
    \item database
    \item tests
\end{enumerate}

That order is much easier than jumping directly into a large auth or repository file without a mental model.

\section{Contact}

\KazeContactBlock

\end{document}
